\subsection{LINEAR LIE ALGEBRAS OF NILPOTENT ENDOMORPHISMS}

Lemma 1. Let \(\mathfrak{n}\) be a Lie subalgebra of \(\mathfrak{gl}(\mathrm{V})\) consisting of nilpotent endomorphisms, and \(\mathrm{N}\) the subgroup \(\exp \mathfrak{n}\) of \(\mathbf{GL}(\mathrm{V})\) (§3, no. 1, Lemma 1).

\begin{enumerate}
\def\labelenumi{(\roman{enumi})}
\item
  Let \(\rho\) be a finite dimensional linear representation of \(\mathfrak{n}\) on W, such that the elements of \(\rho(\mathfrak{n})\) are nilpotent, W' a vector subspace of W stable under \(\rho\) , \(\rho_1\) and \(\rho_2\) the subrepresentation and quotient representation of \(\rho\) defined by W', \(\pi, \pi_1, \pi_2\) the representations of N compatible with \(\rho, \rho_1, \rho_2\) (§3, no. 1). Then \(\pi_1, \pi_2\) are the subrepresentation and quotient representation of \(\pi\) defined by W'.
\item
  Let \(\rho_1, \rho_2\) be finite dimensional linear representations of \(\mathfrak{n}\) such that the elements of \(\rho_1(\mathfrak{n})\) and \(\rho_2(\mathfrak{n})\) are nilpotent, and \(\pi_1, \pi_2\) the representations of N compatible with \(\rho_1, \rho_2\) . Then \(\pi_1 \otimes \pi_2\) is the representation of N compatible with \(\rho_1 \otimes \rho_2\) .
\item
  Let \(\rho_{1},\rho_{2}\) be finite dimensional linear representations of \(\mathfrak{n}\) on vector spaces \(V_{1},V_{2}\) , such that the elements of \(\rho_1(\mathfrak{n})\) and \(\rho_{2}(\mathfrak{n})\) are nilpotent, \(\rho\) the representation of \(\mathfrak{n}\) on \(\mathrm{Hom}(V_1,V_2)\) determined by \(\rho_{1},\rho_{2}\) . Let \(\pi_1,\pi_2\) be the representations of N compatible with \(\rho_{1},\rho_{2}\) , and \(\pi\) the representation of N on \(\mathrm{Hom}(\mathrm{V}_1,\mathrm{V}_2)\) determined by \(\pi_1,\pi_2\) . Then \(\pi\) is the representation of N compatible with \(\rho\) .
\end{enumerate}

Assertion (i) is clear. Let \(\rho_1, \rho_2, \pi_1, \pi_2\) be as in (ii). If \(x \in \mathfrak{n}\) , we have, since \(\rho_1(x) \otimes 1\) and \(1 \otimes \rho_2(x)\) commute,

\[
\begin{array}{l} \exp (\rho_ {1} (x) \otimes 1 + 1 \otimes \rho_ {2} (x)) = \exp (\rho_ {1} (x) \otimes 1). \exp (1 \otimes \rho_ {2} (x)) \\ \qquad = (\exp \rho_ {1} (x)) \otimes 1. 1 \otimes (\exp \rho_ {2} (x)) \\ \qquad = (\exp \rho_ {1} (x)) \otimes (\exp \rho_ {2} (x)) \\ \qquad = \pi_ {1} (\exp x) \otimes \pi_ {2} (\exp x) \\ \qquad = (\pi_ {1} \otimes \pi_ {2}) (\exp x), \end{array}
\]

hence (ii). Let \(\rho_1, \rho_2, \rho, \pi_1, \pi_2, \pi, V_1, V_2\) be as in (iii). If \(v_1 \in \mathrm{End}V_1\) and \(v_2 \in \mathrm{End}V_2\) , denote by \(R_{v_1}\) and \(L_{v_2}\) the maps \(u \mapsto uv_1\) and \(u \mapsto v_2u\) from \(\mathrm{Hom}(V_1, V_2)\) to itself; these maps commute and \(\rho(x)u = (L_{\rho_2(x)} - R_{\rho_1(x)})u\) , so

\[
\begin{array}{r l} & {\exp \rho (x). u = \exp \mathrm{L} _ {\rho_ {2} (x)}. \exp \mathrm{R} _ {- \rho_ {1} (x)}. u} \\ & {\qquad = \mathrm{L} _ {\exp \rho_ {2} (x)}. \mathrm{R} _ {\exp (- \rho_ {1} (x))}. u} \\ & {\qquad = \mathrm{L} _ {\pi_ {2} (\exp x)}. \mathrm{R} _ {\pi_ {1} (\exp (- x))}. u} \\ & {\qquad = \pi (\exp x). u,} \end{array}
\]

hence (iii).

Lemma 2 \(^{2}\) . (i) Let W be a vector subspace of V of dimension d, D the line \(\bigwedge^{d}\) W \(\subset \bigwedge^{d}\) V, \(\theta\) the canonical representation of \(\mathfrak{gl}(V)\) on \(\bigwedge V\) (Chap. III, App.). Let \(x \in \mathfrak{gl}(V)\) . Then \(x(W) \subset W\) if and only if \(\theta(x)(D) \subset D\) .

\begin{enumerate}
\def\labelenumi{(\roman{enumi})}
\setcounter{enumi}{1}
\tightlist
\item
  Let \((e_{1},\ldots,e_{n})\) be the canonical basis of \(k^{n}\) , \(\theta\) the canonical representation of \(\mathfrak{gl}(n,k)\) on \(\bigwedge(k^{n})\) , and \(x\in\mathfrak{gl}(n,k)\) . Then \(x\in\mathfrak{n}(n,k)\) if and only if
\end{enumerate}

\[
\theta (x) (e _ {n - d + 1} \wedge \dots \wedge e _ {n}) = 0
\]

for \(1 \leq d \leq n\) .

\footnotetext[2]{In this lemma, \(k\) can be an arbitrary (commutative) field.}

\begin{enumerate}
\def\labelenumi{(\roman{enumi})}
\tightlist
\item
  If \(x(\mathrm{W}) \subset \mathrm{W}\) , it is clear that \(\theta(x)\mathrm{D} \subset \mathrm{D}\) . Conversely, assume that \(\theta(x)\mathrm{D} \subset \mathrm{D}\) . Let \(u\) be a non-zero element of \(\mathrm{D}\) and let \(y \in \mathrm{W}\) . Then \(y \wedge u = 0\) . Since \(\theta(x)\) is a derivation of \(\bigwedge \mathrm{V}\) , this implies
\end{enumerate}

\[
\theta (x) y \wedge u + y \wedge \theta (x) u = 0.
\]

Now \(\theta(x)u\in ku\) , so \(y\wedge\theta(x)u=0\) and consequently \(\theta(x)y\wedge u=0\) . By Algebra, Chap. III, §7, no. 9, Prop. 13, this implies that \(\theta(x)y\in W\) , i.e.~\(x(y)\in W\) , which proves that \(x(W)\subset W\) .

\begin{enumerate}
\def\labelenumi{(\roman{enumi})}
\setcounter{enumi}{1}
\tightlist
\item
  The condition stated in (ii) is clearly necessary for \(x \in \mathfrak{n}(n, k)\) . Assume that it is satisfied. By (i), \(x\) leaves
\end{enumerate}

\[
k e _ {n - d + 1} + \dots + k e _ {n}
\]

stable, and since this holds for \(d = 1, \ldots, n\) , \(x\) is lower triangular. Put

\[
x = (x _ {i j}) _ {1 \leq i, j \leq n}.
\]

We have \(0 = x(e_n) = x_{nn}e_n\) , so \(x_{nn} = 0\) . Let \(i < n\) , and assume that we have proved that \(x_{jj} = 0\) for \(j > i\) . Then

\[
0 = \theta (x) \left(e _ {i} \wedge e _ {i + 1} \wedge \dots \wedge e _ {n}\right) = x _ {i i} \left(e _ {i} \wedge e _ {i + 1} \wedge \dots \wedge e _ {n}\right),
\]

so \(x_{ii} = 0\) . Thus, \(x \in \mathfrak{n}(n,k)\) .

PROPOSITION 8. Let \(\mathfrak{n}\) be a Lie subalgebra of \(\mathfrak{gl}(\mathrm{V})\) consisting of nilpotent elements, \(\mathfrak{q}\) the normalizer of \(\mathfrak{n}\) in \(\mathfrak{gl}(\mathrm{V})\) . There exists a finite dimensional vector space E, a representation \(\rho\) of \(\mathfrak{gl}(\mathrm{V})\) on E, and a vector subspace F of E, satisfying the following conditions:

\begin{enumerate}
\def\labelenumi{(\roman{enumi})}
\item
  the image under \(\rho\) of a homothety of \(V\) is diagonalizable;
\item
  F is stable under \(\rho(\mathfrak{q})\) ;
\item
  \(\mathfrak{n}\) is the set of \(x\in \mathfrak{gl}(\mathrm{V})\) such that \(\rho (x)(\mathrm{F}) = 0\)
\end{enumerate}

Let \(n = \dim V\) . By Engel's theorem, \(V\) can be identified with \(k^n\) in such a way that \(\mathfrak{n} \subset \mathfrak{n}(n,k)\) . Let \(P\) be the algebra of polynomial functions on \(\mathfrak{gl}(n,k)\) . For \(i = 0,1,\ldots\) , let \(P_i\) be the set of elements of \(P\) homogeneous of degree \(i\) . Let \(N = \exp \mathfrak{n}\) , which is a subgroup of the strictly lower triangular group \(T\) . Let \(J\) be the set of elements of \(P\) that are zero on \(N\) ; this is an ideal in \(P\) . Let \(N_J\) be the set of \(x \in \mathfrak{gl}(n,k)\) such that \(p(x) = 0\) for all \(p \in J\) . Then \(N \subset N_J\) . Conversely, let \(x \in N_J\) . Denote by \(p_{ij}\) the polynomial functions giving the entries of an element of \(\mathfrak{gl}(n,k)\) . The ideal \(J\) contains the \(p_{ij}\) (for \(i < j\) ) and the \(p_{ii} - 1\) ; hence \(x \in T\) . On the other hand, if \(u\) is a linear form on \(\mathfrak{gl}(n,k)\) which is zero on \(\mathfrak{n}\) , there exists \(p_u \in P\) such that \(p_u(z) = u(\log z)\) for all \(z \in T\) ( \(\S 3\) , no. 1, Lemma 1 (i)); we have \(p_u \in J\) , so \(u(\log x) = 0\) . It follows that \(\log x\) belongs to \(\mathfrak{n}\) , so \(x \in N\) , proving that \(N = N_J\) .

For all \(p \in P\) and \(g \in \mathbf{GL}_{n}(k)\) , let \(\lambda(g)p\) be the function \(x \mapsto p(g^{-1}x)\) on \(\mathfrak{gl}(n,k)\) ; then \(\lambda(g)p \in \mathrm{P}\) , \(\lambda(g)\) is an automorphism of the algebra P, and \(\lambda\) is a representation of \(\mathbf{GL}_{n}(k)\) on P which leaves each \(P_{i}\) stable. We show that

\[
\mathrm{N} = \{x \in \mathbf {G L} _ {n} (k) \mid \lambda (x) \mathrm{J} = \mathrm{J} \}.\tag{1}
\]

If \(x \in \mathrm{N}, p \in \mathrm{J}, y \in \mathrm{N}\) , then \((\lambda(x)p)(y) = p(x^{-1}y) = 0\) since \(x^{-1}y \in \mathrm{N}\) ; thus \(\lambda(x)p \in \mathrm{J}\) , so \(\lambda(x)\mathrm{J} = \mathrm{J}\) . Let \(x \in \mathbf{GL}_n(k)\) be such that \(\lambda(x)\mathrm{J} = \mathrm{J}\) ; let \(p \in \mathrm{J}\) ; then \(p(x^{-1}) = (\lambda(x)p)(e) = 0\) , so \(x^{-1} \in \mathrm{N}_{\mathrm{J}} = \mathrm{N}\) and \(x \in \mathrm{N}\) . This proves (i).

The ideal J is of finite type (Commutative Algebra, Chap. III, §2, no. 10, Cor. 2 of Th. 2). Hence, there exists an integer q such that, if \(W = P_{0} + P_{1} + \cdots + P_{q}\) , then \(J \cap W\) generates J as an ideal. Denote by \(\lambda_{j}\) (resp. \(\lambda'\) ) the subrepresentation of \(\lambda\) defined by \(P_{J}\) (resp. by W). By (1),

\[
\mathrm{N} = \{x \in \mathbf {G L} _ {n} (k) \mid \lambda^ {\prime} (x) (\mathrm{J} \cap \mathrm{W}) = \mathrm{J} \cap \mathrm{W} \}.\tag{2}
\]

We show that, for all \(j\) , there exists a representation \(\sigma_{j}\) of the Lie algebra \(\mathfrak{gl}(n,k)\) on \(\mathrm{P}_j\) such that:

\[
\sigma_ {j} | \mathfrak {n} (n, k) \text {   is   compatible   } (\S 3, \text {   no.   } 1) \text {   with   } \lambda_ {j} | T.\tag{3}
\]

\[
\text { For   all } x \in k. 1 _ {n}, \sigma_ {j} (x) \text { is   a   homothety. }\tag{4}
\]

Since \(\lambda_{j}\) is the jth symmetric power of \(\lambda_{1}\) , it suffices to prove the existence of \(\sigma_{1}\) , cf.~Lemma 1. Now \(\lambda_{1}\) is the contragredient representation of the representation \(\gamma\) of \(\mathbf{GL}_{n}(k)\) on \(\mathfrak{gl}(n,k)\) given by

\[
\gamma (x) y = x y, \quad x \in \mathbf {G L} _ {n} (k), \quad y \in \mathfrak {g l} (n, k).
\]

Let \(c\) be the representation of the Lie algebra \(\mathfrak{gl}(n,k)\) on \(\mathfrak{gl}(n,k)\) given by

\[
c (x) y = x y, \quad x, y \in \mathfrak {g l} (n, k).
\]

It is immediate that \(c|\mathfrak{n}(n,k)\) and \(\gamma|T\) are compatible, and that \(c(x)\) is a homothety for all \(x \in k.1_{n}\) . Thus, it suffices to take for \(\sigma_{1}\) the dual representation of c (Chap. I, §3, no. 3).

Now let \(\sigma'\) be the representation of \(\mathfrak{gl}(n,k)\) on W given by the direct sum of the \(\sigma_j\) , \(0 \leq j \leq q\) . In view of (2) and the relations

\[
\lambda^ {\prime} (\exp (x)) = \exp (\sigma^ {\prime} (x)) \text { and } \sigma^ {\prime} (\log (y)) = \log (\lambda^ {\prime} (y)), x \in \mathfrak {n} (n, k), y \in \mathrm{T},
\]

we have

\[
\mathfrak {n} = \{x \in \mathfrak {n} (n, k) \mid \sigma^ {\prime} (x) (\mathrm{J} \cap \mathrm{W}) \subset \mathrm{J} \cap \mathrm{W} \}.\tag{5}
\]

Let \(d = \dim(\mathrm{J} \cap \mathrm{W})\) , and let \(\tau = \bigwedge^{d} \sigma'\) . Let \(\mathrm{D} = \bigwedge^{d} (\mathrm{J} \cap \mathrm{W})\) . By (5) and Lemma 2 (i),

\[
\mathfrak {n} = \{x \in \mathfrak {n} (n, k) \mid \tau (x) (\mathrm{D}) \subset \mathrm{D} \}.\tag{6}
\]

But \(\tau (\mathfrak{n}(n,k))\) consists of nilpotent endomorphisms, so (6) can also be written

\[
\mathfrak {n} = \{x \in \mathfrak {n} (n, k) \mid \tau (x) (\mathrm{D}) = 0 \}.\tag{7}
\]

Now let \(E = \bigwedge^{d} W \oplus \bigwedge^{1} V \oplus \bigwedge^{2} V \oplus \cdots \oplus \bigwedge^{n} V\) ; let \(\rho\) be the direct sum of \(\tau\) and the canonical representations of \(\mathfrak{gl}(n, k)\) on \(\bigwedge^{1} V, \ldots, \bigwedge^{n} V\) . Let \(E_{0} \subset E\) be the sum of \(D = \bigwedge^{d} (J \cap W)\) and the lines generated by \(e_{n-j+1} \wedge \cdots \wedge e_{n}\) for \(j = 1, \ldots, n\) . By (7) and Lemma 2 (ii),

\[
\mathfrak {n} = \{x \in \mathfrak {g l} (\mathrm{V}) \mid \rho (x) (\mathrm{E} _ {0}) = 0 \}.\tag{8}
\]

It is immediate that, if \(x \in k.1_{n}\) , \(\rho(x)\) is diagonalizable. Finally, if F is the set of elements of E annihilated by \(\rho(\mathfrak{n})\) , F is stable under \(\rho(\mathfrak{q})\) (Chap. I, §3, no. 5, Prop. 5), and by (8),

\[
\mathfrak {n} = \{x \in \mathfrak {g l} (\mathrm{V}) \mid \rho (x) (\mathrm{F}) = 0 \}.\tag{9}
\]

